% GENERATED FILE. Edit canonical lessons, then run npm run generate:books.
% canonical-source-hash: fnv1a64:826bb1835c775955
% canonical-lessons: FA-C121-gavmish, FA-C121-susk, FA-C121-kerm, FA-C121-hashare, FA-C121-zarafe

\chapter{Buffalo, Beetle, Worm, Insect, Giraffe}
\label{ch:fa-buffalo-beetle-worm-insect-giraffe}
\hlchaptermodality{121}

\begin{chapteropening}
\textbf{By the end of this chapter:} \emph{I can say a buffalo, a beetle, a worm, an insect and a giraffe.}

Complete the last lesson of chapter 121: I can say a buffalo, a beetle, a worm, an insect and a giraffe.
\end{chapteropening}

\section[\texorpdfstring{\fa{گاومیش} (gâvmish)}{گاومیش (gâvmish)}]{\fa{گاومیش} (gâvmish) — a buffalo}

\label{lesson:FA-C121-gavmish}



\begin{quote}
Before the new one: say the Persian for a spring, then the Persian for a salt desert.
\end{quote}

\subsection*{You'll want to know: \fa{گاومیش}}
\textbf{\fa{گاومیش}} — \emph{gâvmish} — \textquotedblleft{}a buffalo\textquotedblright{}.

\begin{cousinweb}
\emph{gâvmish} is usually explained as \emph{gâv}, \textquotedblleft{}cow, ox\textquotedblright{} (inherited, Middle Persian \emph{gāw}), with \emph{mish}, \textquotedblleft{}ewe, sheep\textquotedblright{}: the buffalo named from two kinds of cattle.
\end{cousinweb}

\begin{grammarlens}[title={What you've built}]
One more word for this chapter.
\end{grammarlens}

\subsection*{Your turn}
\begin{itemize}
\raggedright
  \item \emph{Say it:} \emph{gâvmish}
  \item \emph{Say it:} \emph{gâvmish}, once more
  \item \emph{Say it:} say \emph{gâvmish}
  \item \emph{Recall:} say the Persian for a spring, then the Persian for a salt desert, then say \emph{gâvmish} again
\end{itemize}

\begin{tcolorbox}[breakable,colback=teal!4,colframe=teal!35!black,title={Before you move on}]
What does \fa{گاومیش} mean? (\textquotedblleft{}A buffalo\textquotedblright{}.) Say it once more.
\end{tcolorbox}
\section[\texorpdfstring{\fa{سوسک} (susk)}{سوسک (susk)}]{\fa{سوسک} (susk) — a beetle}

\label{lesson:FA-C121-susk}



\begin{quote}
Before the new one: say the Persian for a salt desert, then the Persian for a buffalo.
\end{quote}

\subsection*{You'll want to know: \fa{سوسک}}
\textbf{\fa{سوسک}} — \emph{susk} — \textquotedblleft{}a beetle\textquotedblright{}.

\begin{cousinweb}
\emph{susk} is the everyday word for a beetle, but no source used here gives its origin, which is not settled.
\end{cousinweb}

\begin{grammarlens}[title={What you've built}]
One more word for this chapter.
\end{grammarlens}

\subsection*{Your turn}
\begin{itemize}
\raggedright
  \item \emph{Say it:} \emph{susk}
  \item \emph{Say it:} \emph{susk}, once more
  \item \emph{Say it:} say \emph{susk}
  \item \emph{Recall:} say the Persian for a salt desert, then the Persian for a buffalo, then say \emph{susk} again
\end{itemize}

\begin{tcolorbox}[breakable,colback=teal!4,colframe=teal!35!black,title={Before you move on}]
What does \fa{سوسک} mean? (\textquotedblleft{}A beetle\textquotedblright{}.) Say it once more.
\end{tcolorbox}
\section[\texorpdfstring{\fa{کرم} (kerm)}{کرم (kerm)}]{\fa{کرم} (kerm) — a worm}

\label{lesson:FA-C121-kerm}



\begin{quote}
Before the new one: say the Persian for a buffalo, then the Persian for a beetle.
\end{quote}

\subsection*{You'll want to know: \fa{کرم}}
\textbf{\fa{کرم}} — \emph{kerm} — \textquotedblleft{}a worm\textquotedblright{}.

\begin{cousinweb}
\emph{kerm} is a native Persian word, inherited rather than borrowed; it is usually traced to Middle Persian \emph{kirm}, \textquotedblleft{}worm\textquotedblright{}. EtymDB takes it back to an Indo-European word for \textquotedblleft{}worm\textquotedblright{}, the word behind Sanskrit \emph{kṛmi}. (It also lists French \emph{crème}, a different word.)
\end{cousinweb}

\begin{grammarlens}[title={What you've built}]
One more word for this chapter.
\end{grammarlens}

\subsection*{Your turn}
\begin{itemize}
\raggedright
  \item \emph{Say it:} \emph{kerm}
  \item \emph{Say it:} \emph{kerm}, once more
  \item \emph{Say it:} say \emph{kerm}
  \item \emph{Recall:} say the Persian for a buffalo, then the Persian for a beetle, then say \emph{kerm} again
\end{itemize}

\begin{tcolorbox}[breakable,colback=teal!4,colframe=teal!35!black,title={Before you move on}]
What does \fa{کرم} mean? (\textquotedblleft{}A worm\textquotedblright{}.) Say it once more.
\end{tcolorbox}
\section[\texorpdfstring{\fa{حشره} (hashare)}{حشره (hashare)}]{\fa{حشره} (hashare) — an insect}

\label{lesson:FA-C121-hashare}



\begin{quote}
Before the new one: say the Persian for a beetle, then the Persian for a worm.
\end{quote}

\subsection*{You'll want to know: \fa{حشره}}
\textbf{\fa{حشره}} — \emph{hashare} — \textquotedblleft{}an insect\textquotedblright{}.

\begin{cousinweb}
\emph{hashare} is borrowed from Arabic \emph{ḥašara}, \textquotedblleft{}insect, small creature\textquotedblright{}. Its Arabic root is \emph{ḥ-š-r} (Lane).
\end{cousinweb}

\begin{grammarlens}[title={What you've built}]
One more word for this chapter.
\end{grammarlens}

\subsection*{Your turn}
\begin{itemize}
\raggedright
  \item \emph{Say it:} \emph{hashare}
  \item \emph{Say it:} \emph{hashare}, once more
  \item \emph{Say it:} say \emph{hashare}
  \item \emph{Recall:} say the Persian for a beetle, then the Persian for a worm, then say \emph{hashare} again
\end{itemize}

\begin{tcolorbox}[breakable,colback=teal!4,colframe=teal!35!black,title={Before you move on}]
What does \fa{حشره} mean? (\textquotedblleft{}An insect\textquotedblright{}.) Say it once more.
\end{tcolorbox}
\section[\texorpdfstring{\fa{زرافه} (zarâfe)}{زرافه (zarâfe)}]{\fa{زرافه} (zarâfe) — a giraffe}

\label{lesson:FA-C121-zarafe}



\begin{quote}
Before the new one: say the Persian for a worm, then the Persian for an insect.
\end{quote}

\subsection*{You'll want to know: \fa{زرافه}}
\textbf{\fa{زرافه}} — \emph{zarâfe} — \textquotedblleft{}a giraffe\textquotedblright{}.

\begin{cousinweb}
\emph{zarâfe} is borrowed from Arabic \emph{zarāfa}, \textquotedblleft{}giraffe\textquotedblright{} (EtymDB), and English \emph{giraffe} comes from the same Arabic word. EtymDB traces the Arabic word itself back to a Persian compound, but that is not settled.
\end{cousinweb}

\begin{grammarlens}[title={What you've built}]
One more word for this chapter.
\end{grammarlens}

\subsection*{Your turn}
\begin{itemize}
\raggedright
  \item \emph{Say it:} \emph{zarâfe}
  \item \emph{Say it:} \emph{zarâfe}, once more
  \item \emph{Say it:} say \emph{zarâfe}
  \item \emph{Recall:} say the Persian for a worm, then the Persian for an insect, then say \emph{zarâfe} again
\end{itemize}

\begin{tcolorbox}[breakable,colback=teal!4,colframe=teal!35!black,title={Before you move on}]
What does \fa{زرافه} mean? (\textquotedblleft{}A giraffe\textquotedblright{}.) Say it once more.
\end{tcolorbox}
